% !PRINTED 667
\weberpage{667}
Die Definition der Ordnungszahl haftet also am Punkte $\mathfrak{P}$. Sie ist für den Körper $\Omega$ invariant.

Daß jede Funktion in jedem Punkte $\mathfrak{P}$ eine bestimmte \gesperrt{ganze Zahl} $n$ als Ordnungszahl erhält, ergibt sich nun aus § 185, (5).

Man nehme eine Funktion $z$, die in $\mathfrak{P}$ endlich bleibt, als unabhängige Variable, bestimme das zu $\mathfrak{P}$ gehörige Primfunktional $\pi$ und wähle eine Funktion $\varrho$, deren Zähler durch $\pi$, aber nicht durch $\pi^2$ teilbar ist. Diese Funktion ist von der ersten Ordnung.

Dann kann man jede Funktion $\omega$ in die Form setzen:

\weq{1}{\omega = a\varrho^n + \eta_1\varrho^{n+1},}
worin $a$ eine von Null verschiedene Konstante ist, und $n$ ist die \gesperrt{Ordnungszahl} von $\omega$.

Hat $\varrho$ die Ordnung 1, so hat $\varrho^n$ die Ordnung $n$.

Über die Ordnungszahlen gelten folgende Sätze, die alle leicht aus (1) folgen:

4. \gesperrt{Die Ordnung eines Produktes zweier Funktionen ist gleich der Summe der Ordnungen der Faktoren. Die Ordnung eines Quotienten ist gleich der Differenz der Ordnungen von Zähler und Nenner.}

\gesperrt{Die Ordnung einer Summe ist gleich der niedrigsten unter den Ordnungen der Summanden. Kommen unter den Summanden mehrere von gleicher niedrigster Ordnungszahl vor, so kann die Ordnung der Summe größer, aber nicht kleiner sein.}

\gesperrt{Von einer Funktion, die in $\mathfrak{P}$ Null in der $n$ten Ordnung wird, sagt man auch, sie wird Null in der ersten Ordnung in $n$ zusammengefallenen Punkten.}

\begin{center}
\textbf{§ 189. \gesperrt{Polygone.}}
\end{center}

1. \gesperrt{Definition. Komplexe von Punkten, die denselben Punkt auch mehrmals enthalten können, heißen Polygone oder Vielecke} (des Körpers $\Omega$).

Die Zahl der Punkte eines Polygons heißt seine \gesperrt{Ordnung}, und ein Polygon $m$ter Ordnung wird auch kurz ein $m$-Eck genannt.

Als Bezeichnung für die Polygone sollen die Buchstaben des großen deutschen Alphabets $\mathfrak{A}, \mathfrak{B}, \mathfrak{C}, \ldots$ dienen.
